\thispagestyle{empty}
\begin{ChapterAbstract}
\end{ChapterAbstract}
%\includepdf[pages={1}]{extern/page_vide.pdf} 

%-----------------------------------------------------------------
%-----------------debut de chapitre--------------------------------
%------------------------------------------------------------------

\newpage
\section{Historique}
Avant l'informatisation, la gestion des bibliothèques reposait entièrement sur des procédés manuels et indépendants : bons de commande sur papier, registres d'inventaire manuscrits, fiches cartonnées pour le catalogage, et suivi des prêts assuré par de simples fiches nominatives au bureau de circulation 

Les années 1960 marquent les premiers pas de l'informatisation avec l'invention des formats MARC \cite{marc}, permettant le stockage et l'échange standardisé des notices catalographiques. La décennie suivante voit émerger les premiers Systèmes de Gestion de Bibliothèque (SGB), d'abord centrés sur la circulation des documents, avant de s'enrichir dans les années 1980 de fonctions intégrées : acquisitions, catalogage, gestion des périodiques, réservations, etc.

Les années 1990 apportent Internet et une multiplication des applications. Puis, à partir des années 2010, la dématérialisation transforme profondément le paysage documentaire — livres électroniques, articles scientifiques en ligne, vidéos et musiques en streaming — avec des ressources désormais souvent hébergées chez des fournisseurs externes, disposant de leurs propres plateformes et interfaces de recherche.

\begin{figure}[h]
    \centering
    \begin{minipage}{0.48\textwidth}
        \centering
        \includegraphics[width=\linewidth]{chapter_1/figures/catalogue.jpeg}
    \end{minipage}
    \hfill
    \begin{minipage}{0.48\textwidth}
        \centering
        \includegraphics[width=\linewidth]{chapter_1/figures/cover.jpeg}
    \end{minipage}
    \caption{Source: Gallica Bnf, «Palais impérial de Compiègne. Bibliothèque. Registre de prêts de livres,1869-1870»}
    \label{fig:registres_anciens}
\end{figure}
% ---------------------------------------------------------------------------
\section{Les systèmes intégrés de gestion de bibliothèque (SIGB)}

\subsection{Définition et rôle}

Un \textbf{Système Intégré de Gestion de Bibliothèque} (SIGB), également désigné
sous l'appellation Système de Gestion de Bibliothèque (SGB), est un progiciel
destiné à l'automatisation des différentes activités nécessaires au fonctionnement
d'une bibliothèque. Il s'appuie sur une base de données relationnelle pour
proposer des modules de gestion couvrant l'ensemble du cycle de vie du document :
acquisition, catalogage, circulation, consultation et administration
\cite{Chowdhury_Burton_McMenemy_Poulter_2007}.

Ces outils visent à libérer le personnel des tâches répétitives — préparation
des commandes, bulletinage, suivi des retards — afin de le recentrer sur des
missions à plus forte valeur ajoutée, notamment l'accueil et le conseil aux
usagers \cite{oyekale2018ils}
% \cite{mataoui2022}
. Parmi les bénéfices concrets d'un SIGB, on peut citer :

\begin{itemize}
\item la traçabilité continue de la disponibilité des exemplaires ;
\item la diminution des erreurs liées aux traitements manuels ;
\item l’accès à une recherche documentaire multicritère via un catalogue en ligne ;
\item la production automatisée de statistiques d’utilisation ;
\item l’automatisation des relances ainsi que du calcul des pénalités de retard.
\end{itemize}

Un SIGB peut être étudié selon deux axes complémentaires : les opérations
orientées \emph{document} (catalogage, indexation, acquisition) et les
opérations orientées \emph{lecteur} (prêt, réservation, pénalités)
% \cite{mataoui2022}
.

\subsection{Modules fonctionnels d'un SIGB}

Un SIGB complet s'articule généralement autour des modules fonctionnels
suivants 
% \cite{mataoui2022} 
:

\begin{itemize}
    \item \textbf{Module d'acquisition} : gestion des suggestions d'achat,
    des commandes, de la réception des documents, des fournisseurs et du budget.
    Il permet de suivre les engagements financiers par type de document ou
    par département.

    \item \textbf{Module de catalogage} : saisie ou import de notices
    bibliographiques au format UNIMARC, gestion des autorités (auteurs,
    éditeurs, collectivités), classification selon la CDD ou la CDU, et
    création des exemplaires physiques.

    \item \textbf{Module des périodiques} : gestion des abonnements,
    bulletinage des numéros reçus, suivi des états de collection et gestion
    des reliures.

    \item \textbf{Module de circulation (prêt)} : gestion des prêts, retours,
    renouvellements et réservations, avec paramétrage des durées et quotas
    par catégorie de lecteur, et calcul automatique des pénalités en cas de
    retard.

    \item \textbf{Module lecteurs (adhérents)} : gestion des fiches adhérents,
    des catégories et des états (actif, suspendu), consultation de l'historique
    des prêts.

    \item \textbf{Module OPAC} (\emph{Online Public Access Catalog}) :
    interface de recherche documentaire publique proposant des recherches
    simple et avancée (multicritères, opérateurs booléens, troncature).

    \item \textbf{Module statistiques et administration} : génération de
    rapports, paramétrage du système, gestion des droits d'accès et
    sauvegarde des données.
\end{itemize}

\subsection{Panorama des solutions existantes et tableau comparatif}
\label{sec:comparatif}

Le marché des SIGB est segmenté entre des solutions \emph{open-source}
librement accessibles et des solutions \emph{propriétaires} commerciales.
Le tableau~\ref{tab:comparatif_sigb} présente un aperçu comparatif des
principales solutions de référence selon deux axes : fonctionnel (modules
couverts) et technique (architecture, technologies, licence).
\begin{table}[H]
\centering
\caption{Comparatif des principaux SIGB open-source et propriétaires}
\label{tab:comparatif_sigb}
\renewcommand{\arraystretch}{1.4}
\small

\begin{tabularx}{\textwidth}{|p{2.1cm}|p{1.9cm}|p{1.5cm}|p{1.9cm}|p{1.5cm}|X|}
\hline
\textbf{Logiciel} & \textbf{Licence} & \textbf{Techno.} & \textbf{SGBD} &
\textbf{Modules} & \textbf{Remarques} \\
\hline
\textbf{Koha} & Opensource & Perl / JS & MySQL / MariaDB &
Complet & Référence mondiale, UNIMARC/MARC21, très actif \\
\hline
\textbf{PMB} & Opensource & PHP & MySQL &
Complet & Solution française, OPAC intégré, communauté active \\
\hline
\textbf{Evergreen} & Opensource & C / JS & PostgreSQL &
Complet & Conçu pour réseaux de bibliothèques \\
\hline
\textbf{OpenBiblio} & Opensource & PHP & MySQL &
Prêt, cat. & Inactif depuis 2014, fonctionnalités limitées \\
\hline
\textbf{Invenio} & Opensource & Python & PostgreSQL &
Cat., dépôt & Orienté archives institutionnelles (CERN) \\
\hline
\textbf{Alma} (Ex Libris) & Propriétaire & Cloud & SaaS &
Complet+ & Standard universitaire mondial, coût élevé \\
\hline
\textbf{SYNGEB} (CERIST) & Propriétaire & Delphi & SQL Server &
Prêt, cat. & Utilisé en Algérie, mono-poste, non maintenu \\
\hline
\textbf{CATLIB} (EMP) & Interne & Delphi / PHP & Oracle 11g &
Prêt, cat., OPAC & Solution sur mesure, 61 tables, intranet EMP \\
\hline
\end{tabularx}
\end{table}

\FloatBarrier
\newpage
Cette comparaison met en évidence une réalité importante : les solutions
\emph{open-source} matures telles que Koha ou PMB proposent des ensembles
fonctionnels complets et maintenus par des communautés actives, ces deux progiciels sont classés parmi les plus choisis par les bibliothèques universitaires en 2025 (figure~\ref{fig:tosca_consultants}) . Cependant,
leur déploiement dans un contexte institutionnel militaire soulève des
contraintes spécifiques : intégration avec un serveur d'authentification
interne, adaptation à un schéma de données existant de 61 tables, et
maîtrise totale de l'infrastructure réseau. Ces contraintes justifient
pleinement le choix d'un développement sur mesure pour la modernisation
du module de prêt.

\begin{figure}[H]
    \centering
    \includegraphics[width=0.75\textwidth]{chapter_1/figures/progiciels-choisis-par-les-bibliotheques-univrsitaires-en-2025-1024x615.png}
    \caption{Tosca Consultants,Résultats de l’enquête 2026}
    \label{fig:tosca_consultants}
\end{figure}

% ---------------------------------------------------------------------------
\section{Contexte de l'établissement}

\subsection{Présentation de l'École Militaire Polytechnique}

L'École Militaire Polytechnique (EMP) est un établissement d'enseignement
supérieur et de recherche scientifique à vocation militaire et technique.
Elle forme des cadres dans les domaines de l'ingénierie et des sciences
appliquées, et dispose à ce titre d'une infrastructure documentaire importante,
gérée par le \textbf{Bureau de l'Information Scientifique et Technique (BIST)}.

Le BIST est organisé autour du Centre de Documentation (CDOC), lui-même
structuré en plusieurs sections : acquisitions, périodiques, prêt et
catalogage. Il compte plus de dix agents et gère un fonds documentaire de
plus de \textbf{50\,000 titres} pour une base d'environ \textbf{1\,500
adhérents} appartenant à différentes catégories (personnel enseignant,
chercheurs, élèves officiers, personnel administratif)
.
% \cite{mataoui2022}.
\newpage
\subsection{Historique de l'informatisation}

L'informatisation du BIST s'est déroulée en plusieurs étapes successives,
résumées dans le tableau~\ref{tab:historique}.

\begin{table}[H]
\centering
\caption{Chronologie de l'informatisation de la bibliothèque de l'EMP}
\label{tab:historique}
\renewcommand{\arraystretch}{1.45}
\begin{tabular}{|p{2.6cm}|p{10.5cm}|}
\hline

\textbf{Période} & \textbf{Étape} \\
\hline
Avant 2000 &
Gestion entièrement manuelle : fiches cartonnées catalographiques, registres
d'inventaire manuscrits, fiches d'emprunt glissées dans les documents,
vérification quotidienne manuelle des retards par les agents. \\
\hline
Début 2000s &
Acquisition du logiciel commercial \textbf{SYNGEB} (CERIST) : solution
mono-poste, à code fermé, avec abonnement et mises à jour payants, pas
d'OPAC, paramétrage non adapté aux besoins spécifiques de l'EMP. \\
\hline
2007 &
Décision de développer un SIGB en interne. Étude de l'existant, définition
des besoins à court, moyen et long terme, choix technologiques (Oracle 11g,
Delphi, PHP). \\
\hline
2007–2011 &
Développement progressif de \textbf{CATLIB} en cinq étapes : extraction et
migration des données depuis SYNGEB, création de la base Oracle (61 tables),
module OPAC web (PHP), module de prêt Delphi, modules d'acquisition et de
catalogage. \\
\hline
2011–présent &
Mise en production et exploitation. 90\,\% du fonds documentaire numérisé.
Déployé sous forme de machines virtuelles sur l'intranet de l'établissement. \\
\hline
2026 &
Lancement du présent projet : modernisation du module de prêt en
\textbf{.NET\,10 / Blazor Server}, en conservant la base Oracle existante. \\
\hline
\end{tabular}
\end{table}

\FloatBarrier

% ---------------------------------------------------------------------------
\section{L'application existante : CATLIB}

\subsection{Architecture technique}

Le système CATLIB repose sur une architecture \textbf{client-serveur} articulée
autour des composantes suivantes
% \cite{mataoui2022} 
:

\begin{itemize}
    \item \textbf{SGBD} : Oracle Database 11g, hébergeant la base de données composée de 61 tables organisées en cinq groupes
    fonctionnels (notices bibliographiques, circulation, indexation, thèses
    et mémoires, informations bibliographiques complémentaires).

    \item \textbf{Application de catalogage} : client lourd développé en
    \textbf{Delphi}, permettant la saisie et la mise à jour des notices,
    l'indexation des termes et la création des exemplaires
    (figure~\ref{fig:catlib_data}).

    \item \textbf{Application de prêt (CATLIB-Prêt)} : client lourd développé
    en \textbf{Delphi 6 / Object Pascal}, couvrant toutes les opérations de
    circulation (figure~\ref{fig:catlib_menu}).

    \item \textbf{Portail OPAC (CATLIB Web)} : application web intranet
    développée en \textbf{PHP/HTML/JS}, proposant la recherche documentaire
    simple et avancée avec consultation en plein texte via Adobe Flash Player
    (figure~\ref{fig:catlib_opac}).
\end{itemize}

La conception de la base de données a été réalisée selon la méthode
\textbf{MERISE}. Le schéma comporte 61 tables regroupées en cinq parties :

\begin{enumerate}
    \item \textbf{Notices bibliographiques} : \texttt{NOTICE},
    \texttt{TYPE\_NOTICE}, \texttt{SOURCE\_ARTICLE}, \texttt{PERIODICITE},
    \texttt{NOTICE\_COLLECTION}, \ldots

    \item \textbf{Circulation et adhérents} : \texttt{PRET},
    \texttt{RESERVATION}, \texttt{HISTORIQUE\_PRET}, \texttt{ADHERENT},
    \texttt{CATEGORIE}, \texttt{EXEMPLAIRE}, \texttt{ETAT\_EXEMPLAIRE},
    \texttt{PENALITE}, \texttt{PENALITE\_ADHERENT},
    \texttt{PENALITE\_ADHERENT\_TEMP}, \texttt{HISTORIQUE\_PENALITE\_ADHERENT},
    \texttt{JOURS\_FERIES}, \texttt{PARAMETRES\_CATLIB\_PRET}, \ldots

    \item \textbf{Indexation textuelle} : \texttt{TERME}, \texttt{TERME\_EXACT},
    \texttt{NOTICE\_TERME}, \texttt{NOTICE\_TERME\_EXACT}, \texttt{MOTS\_CLES},
    \texttt{MOTS\_VIDES}, \ldots

    \item \textbf{Thèses et mémoires} : \texttt{NOTICE\_DIP\_DIS\_ETAB},
    \texttt{DIPLOME}, \texttt{DISCIPLINE}, \texttt{ETABLISSEMENT}, \ldots

    \item \textbf{Informations bibliographiques} : \texttt{EDITEUR},
    \texttt{VILLE}, \texttt{PAYS}, \texttt{LANGUE},
    \texttt{MENTION\_RESPONSABILITE}, \texttt{AUTEUR},
    \texttt{AUTEUR\_SECONDAIRE}, \texttt{CO\_AUTEUR}, \ldots
\end{enumerate}

\subsection{Description des écrans de CATLIB-Prêt}

L'application CATLIB-Prêt est l'interface de travail quotidienne des agents
du BIST pour la gestion de la circulation. Son menu principal
(figure~\ref{fig:catlib_menu}) donne accès à l'ensemble des fonctions
de circulation.
\begin{figure}[H]
    \centering
    \includegraphics[width=0.75\textwidth]{chapter_1/figures/11.png}
    \caption{Menu principal de CATLIB-Prêt (application Delphi)}
    \label{fig:catlib_menu}
\end{figure}

\subsubsection{Prêt}

Le formulaire de prêt (figure~\ref{fig:catlib_pret}) permet à l'agent de
saisir le numéro de carte adhérent et la cote du document. L'identité de
l'adhérent est automatiquement affichée après reconnaissance. Un menu
déroulant liste les exemplaires disponibles correspondant à la cote saisie.
La date de prêt est pré-remplie avec la date courante ; la date de retour
prévue est calculée automatiquement selon la durée de prêt associée à la
catégorie de l'adhérent (champ \texttt{DUREE\_PRET} de la table
\texttt{CATEGORIE}).
\begin{figure}[H]
    \centering
    \includegraphics[width=0.75\textwidth]{chapter_1/figures/12.png}
    \caption{Formulaire de saisie d'un prêt dans CATLIB-Prêt}
    \label{fig:catlib_pret}
\end{figure}
\subsubsection{Restitution et renouvellement}

L'écran de restitution (figure~\ref{fig:catlib_restitution}) permet de
valider le retour d'un exemplaire. Il affiche les dates de prêt et de
retour effectif, et propose également un bouton \textbf{Renouvellement}
pour prolonger la durée de prêt d'un document non encore rendu.
\begin{figure}[H]
    \centering
    \includegraphics[width=0.75\textwidth]{chapter_1/figures/13.png}
    \caption{Formulaire de restitution et renouvellement}
    \label{fig:catlib_restitution}
\end{figure}

\subsubsection{Avis de disponibilité}

L'écran d'avis de disponibilité (figure~\ref{fig:catlib_dispo}) liste
les adhérents ayant réservé un document qui vient d'être retourné.
L'agent peut envoyer les notifications par courrier électronique ou
les imprimer.
\begin{figure}[H]
    \centering
    \includegraphics[width=0.72\textwidth]{chapter_1/figures/14.png}
    \caption{Écran des avis de disponibilité (réservations)}
    \label{fig:catlib_dispo}
\end{figure}

\subsubsection{Gestion des adhérents}

Le formulaire de gestion des adhérents (figure~\ref{fig:catlib_adherents})
permet la recherche par identifiant ou par nom, la consultation et la
modification des informations de l'adhérent (nom, prénom, position,
catégorie, état).
\begin{figure}[H]
    \centering
    \includegraphics[width=0.72\textwidth]{chapter_1/figures/16.png}
    \caption{Formulaire de gestion des adhérents}
    \label{fig:catlib_adherents}
\end{figure}

\subsubsection{Gestion des documents}

L'écran de gestion des documents (figure~\ref{fig:catlib_docs}) offre une
vue consolidée par exemplaire : prêt en cours, réservation active et
historique des prêts passés.
\begin{figure}[H]
    \centering
    \includegraphics[width=0.80\textwidth]{chapter_1/figures/17.png}
    \caption{Écran de gestion des documents (état d'un exemplaire par cote)}
    \label{fig:catlib_docs}
\end{figure}

\FloatBarrier

% ---------------------------------------------------------------------------
\section{Critique de l'existant et cahier des charges}

\subsection{Limitations techniques}

\subsubsection{Obsolescence de Delphi 6 et d'Adobe Flash}

L'application CATLIB-Prêt est développée en \textbf{Delphi 6}, environnement
de développement publié en 2001 dont le support commercial actif a cessé.
Le compilateur Object Pascal 32 bits sous-jacent présente des incompatibilités
avec les systèmes Windows 64 bits modernes, nécessitant des couches de
compatibilité WoW64, une machine virtuelle ou des configurations spécifiques pour fonctionner.

\newpage
Par ailleurs, l'OPAC CATLIB Web utilise \textbf{Adobe Flash Player} pour
la consultation en plein texte des documents numérisés. Flash a été
définitivement abandonné par Adobe en décembre 2020 et est bloqué nativement
par l'ensemble des navigateurs modernes. Cette fonctionnalité est donc
actuellement non opérationnelle.

\subsubsection{Architecture client lourd}

CATLIB-Prêt est un \textbf{exécutable Win32} devant être installé localement
sur chaque poste de travail. Cette contrainte implique une gestion du
déploiement poste par poste, l'impossibilité d'un accès concurrent multi-poste
géré côté application, et une dépendance aux bibliothèques runtime Delphi
présentes sur la machine cliente.

\subsubsection{Gestion de l'authentification insuffisante}

10:45 PML'examen du schéma de la table ADMIN révèle que l'authentification est uniquement implémentée pour les administrateurs, les mots de passe étant stockés en clair sans aucun mécanisme de protection. Aucune authentification n'est prévue pour les autres utilisateurs du système, et il n'existe aucune intégration avec l'annuaire institutionnel de l'EMP, obligeant à une gestion d'identités distincte et redondante.

\subsection{Limitations fonctionnelles et ergonomiques}

\subsubsection{Interface utilisateur datée}

Les formulaires Delphi présentent une interface graphique figée, sans retour
visuel immédiat sur les erreurs de saisie. Il n'existe pas de tableau de bord
synthétique permettant de visualiser l'activité de la bibliothèque en temps
réel (emprunts du jour, retards en cours, réservations à notifier).

\subsection{Récapitulatif des problèmes identifiés}

\begin{table}[H]
\centering
\caption{Récapitulatif des limitations du système CATLIB existant}
\label{tab:critique}
\renewcommand{\arraystretch}{1.45}
\begin{tabular}{|p{2.8cm}|p{8.5cm}|p{2cm}|}
\hline

\textbf{Dimension} & \textbf{Problème identifié} & \textbf{Sévérité} \\
\hline
Technique & Delphi 6 obsolète, incompatible 64 bits & Critique \\
\hline
Technique & Adobe Flash supprimé, OPAC plein texte non fonctionnel & Critique \\
\hline
Technique & Client lourd, déploiement poste par poste & Élevée \\
\hline
Sécurité & Mots de passe stockés en clair dans Oracle & Critique \\
\hline
Sécurité & Aucune intégration avec l'annuaire institutionnel & Élevée \\
\hline

Ergonomie & Interface datée, absence de tableau de bord & Modérée \\
\hline

\end{tabular}
\end{table}

\FloatBarrier

\subsection{Cahier des charges}

\subsubsection{Besoins fonctionnels}

\begin{itemize}
    \item \textbf{BF-01} — Gestion complète du cycle de prêt : enregistrement,
    consultation et clôture d'un prêt, avec calcul automatique de la date de
    retour prévue selon la catégorie de l'adhérent.
    \item \textbf{BF-02} — Restitution avec détection automatique des retards
    et application immédiate des pénalités.
    \item \textbf{BF-03} — Renouvellement d'un prêt en cours, sous réserve
    d'absence de pénalité active.
    \item \textbf{BF-04} — Gestion des réservations : création, notification
    de disponibilité et expiration automatique après délai paramétrable.
    \item \textbf{BF-05} — Mise à jour quotidienne automatisée des pénalités
    et des réservations expirées via un programme console planifié (tâche
    Windows Scheduler).
    \item \textbf{BF-06} — Tableau de bord en temps réel : emprunts du jour,
    retards en cours, réservations en attente.
    \item \textbf{BF-07} — Fiche adhérent avec historique complet des prêts,
    pénalités et réservations.
\end{itemize}

\subsubsection{Besoins non-fonctionnels}

\begin{itemize}
    \item \textbf{BNF-01} — \textbf{Accessibilité web} : accès depuis tout
    poste du réseau intranet, sans installation locale.
    \item \textbf{BNF-02} — \textbf{Authentification centralisée} : intégration
    avec le serveur d'authentification institutionnel de l'EMP.
    \item \textbf{BNF-03} — \textbf{Interopérabilité} : conservation et
    exploitation de la base Oracle existante (61 tables), sans migration
    de schéma.
    \item \textbf{BNF-04} — \textbf{Maintenabilité} : architecture en couches
    (Controller / Service / Repository), injection de dépendances, couverture
    par tests unitaires et d'intégration.
    \item \textbf{BNF-05} — \textbf{Performance} : temps de réponse des
    opérations courantes (prêt, retour) inférieur à deux secondes.
    \item \textbf{BNF-06} — \textbf{Sécurité} : aucun stockage de mot de passe
    dans l'application, gestion des rôles (agent, administrateur).
\end{itemize}

\subsubsection{Périmètre du projet}

Ce projet couvre le \textbf{module de circulation} du SIGB : prêt,
restitution, réservation, renouvellement, pénalités, ainsi que le programme
de traitement automatisé quotidien associé. Les modules de catalogage,
d'acquisition et l'OPAC, déjà opérationnels dans le cadre de CATLIB,
sont hors périmètre mais l'architecture retenue est conçue pour accueillir
leur modernisation future.

% ---------------------------------------------------------------------------
\section*{Conclusion du chapitre}
\addcontentsline{toc}{section}{Conclusion du chapitre }

Ce premier chapitre a établi le cadre du projet à travers trois axes.
Après une présentation des SIGB et un état de l'art des solutions disponibles
sur le marché, nous avons retracé l'historique de l'informatisation du BIST
de l'EMP et décrit en détail le système CATLIB développé entre 2007 et 2011.
L'analyse critique de cette solution a mis en évidence des limitations
techniques majeures — obsolescence de Delphi 6, disparition d'Adobe Flash,
architecture client lourd — ainsi que des manques fonctionnels et de sécurité
qui motivent une refonte du module de prêt. Le cahier des charges qui en découle
constitue la base des chapitres suivants, consacrés à la conception puis à la
réalisation de la nouvelle application.